\documentclass[10pt,a4paper]{amsart}
\usepackage{amsmath,amssymb,amsthm,graphicx}
\usepackage{verbatim}
\newenvironment{reference}{\comment}{\endcomment}
%\newenvironment{reference}{}{}
\newenvironment{slogan}{\comment}{\endcomment}
\newenvironment{history}{\comment}{\endcomment}

% For commutative diagrams we use Xy-pic
\usepackage[all]{xy}

% We use 2cell for 2-commutative diagrams.
\xyoption{2cell}
% Theorem environments.
%
\theoremstyle{plain}
\newtheorem{theorem}[subsection]{Theorem}
\newtheorem{proposition}[subsection]{Proposition}
\newtheorem{lemma}[subsection]{Lemma}

\theoremstyle{definition}
\newtheorem{definition}[subsection]{Definition}
\newtheorem{example}[subsection]{Example}
\newtheorem{exercise}[subsection]{Exercise}
\newtheorem{situation}[subsection]{Situation}

\theoremstyle{remark}
\newtheorem{remark}[subsection]{Remark}
\newtheorem{remarks}[subsection]{Remarks}

\numberwithin{equation}{subsection}

% Macros
%
\def\lim{\mathop{\mathrm{lim}}\nolimits}
\def\colim{\mathop{\mathrm{colim}}\nolimits}
\def\Spec{\mathop{\mathrm{Spec}}}
\def\Hom{\mathop{\mathrm{Hom}}\nolimits}
\def\Ext{\mathop{\mathrm{Ext}}\nolimits}
\def\SheafHom{\mathop{\mathcal{H}\!\mathit{om}}\nolimits}
\def\SheafExt{\mathop{\mathcal{E}\!\mathit{xt}}\nolimits}
\def\Sch{\mathit{Sch}}
\def\Mor{\mathop{\mathrm{Mor}}\nolimits}
\def\Ob{\mathop{\mathrm{Ob}}\nolimits}
\def\Sh{\mathop{\mathit{Sh}}\nolimits}
\def\NL{\mathop{N\!L}\nolimits}
\def\CH{\mathop{\mathrm{CH}}\nolimits}
\def\proetale{{pro\text{-}\acute{e}tale}}
\def\etale{{\acute{e}tale}}
\def\QCoh{\mathit{QCoh}}
\def\Ker{\mathop{\mathrm{Ker}}}
\def\Im{\mathop{\mathrm{Im}}}
\def\Coker{\mathop{\mathrm{Coker}}}
\def\Coim{\mathop{\mathrm{Coim}}}

% Boxtimes
%
\DeclareMathSymbol{\boxtimes}{\mathbin}{AMSa}{"02}

%
% Macros for moduli stacks/spaces
%
\def\QCohstack{\mathcal{QC}\!\mathit{oh}}
\def\Cohstack{\mathcal{C}\!\mathit{oh}}
\def\Spacesstack{\mathcal{S}\!\mathit{paces}}
\def\Quotfunctor{\mathrm{Quot}}
\def\Hilbfunctor{\mathrm{Hilb}}
\def\Curvesstack{\mathcal{C}\!\mathit{urves}}
\def\Polarizedstack{\mathcal{P}\!\mathit{olarized}}
\def\Complexesstack{\mathcal{C}\!\mathit{omplexes}}
% \Pic is the operator that assigns to X its picard group, usage \Pic(X)
% \Picardstack_{X/B} denotes the Picard stack of X over B
% \Picardfunctor_{X/B} denotes the Picard functor of X over B
\def\Pic{\mathop{\mathrm{Pic}}\nolimits}
\def\Picardstack{\mathcal{P}\!\mathit{ic}}
\def\Picardfunctor{\mathrm{Pic}}
\def\Deformationcategory{\mathcal{D}\!\mathit{ef}}
\usepackage{cleveref}
\crefname{originallemma}{original lemma}{original lemmas}
\Crefname{originallemma}{Original lemma}{Original lemmas}
\makeatletter
\newlabel{algebra-lemma-dim-affine-space}{{00OP}{source}{}{sourcepointer-1}{}}
\newlabel{algebra-lemma-dim-affine-space@cref}{{[originallemma][0][]00OP}{[1][0][]source}{}{}{}}
\makeatother
\crefname{originallemma}{original lemma}{original lemmas}
\Crefname{originallemma}{Original lemma}{Original lemmas}
\makeatletter
\newlabel{algebra-lemma-regular-domain}{{00NP}{source}{}{sourcepointer-2}{}}
\newlabel{algebra-lemma-regular-domain@cref}{{[originallemma][0][]00NP}{[1][0][]source}{}{}{}}
\makeatother
\crefname{originallemma}{original lemma}{original lemmas}
\Crefname{originallemma}{Original lemma}{Original lemmas}
\makeatletter
\newlabel{algebra-lemma-regular-mod-x}{{00NU}{source}{}{sourcepointer-3}{}}
\newlabel{algebra-lemma-regular-mod-x@cref}{{[originallemma][0][]00NU}{[1][0][]source}{}{}{}}
\makeatother
\crefname{originallemma}{original lemma}{original lemmas}
\Crefname{originallemma}{Original lemma}{Original lemmas}
\makeatletter
\newlabel{divisors-lemma-blow-up-integral-scheme}{{02ND}{source}{}{sourcepointer-4}{}}
\newlabel{divisors-lemma-blow-up-integral-scheme@cref}{{[originallemma][0][]02ND}{[1][0][]source}{}{}{}}
\makeatother
\crefname{originallemma}{original lemma}{original lemmas}
\Crefname{originallemma}{Original lemma}{Original lemmas}
\makeatletter
\newlabel{properties-lemma-characterize-regular}{{02IT}{source}{}{sourcepointer-5}{}}
\newlabel{properties-lemma-characterize-regular@cref}{{[originallemma][0][]02IT}{[1][0][]source}{}{}{}}
\makeatother
\title{Regularity of the blowup of the maximal ideal of a regular local surface}
\author{The Stacks Project Authors}
\date{stacks-project:resolve.tex:local-snapshot-20261012; GNU FDL 1.2 or later}
\begin{document}
\maketitle
\noindent\textit{Editorial note.} All delivered mathematical text is the author's own, copied from the pinned Stacks Project chapter file resolve.tex; no statement or proof step was rewritten, reordered or removed.\par\medskip
\noindent\textit{Editorial note.} The counted proof is the author's complete argument: the reduction to closed points, the local equation of the exceptional divisor at a closed point, and the quotient criterion for regularity.\par\medskip
\noindent\textit{Editorial note.} The label lemma-blowup-regular names the counted statement, and the proof environment immediately following it in the source is delivered verbatim as its proof.\par\medskip
\noindent\textit{Editorial note.} The counted statement is carried by the author's own lemma environment in the chapter source.\par\medskip
\noindent\textit{Editorial note.} Pre-open gate: python operations/agent-authoring/tools/precard\_gate.py --source resolve.tex --label lemma-blowup-regular --start 726 --end 753 --local-id 28106 returned RESULT=OK (L28106 free, occupancy FREE, public interval clear). Pre-assembly tag self-check: canonical tag resolved and not present in reference/public-items.jsonl.\par\medskip
\noindent\textit{Editorial note.} The statement and proof are copied from the pinned chapter file; the cited inputs are kept as external pointers with their canonical Stacks tags.\par\medskip
\noindent\textit{Allowed standard prerequisites (editorial).}\begin{itemize}
\item Integrality of the blowup of an integral ring in an ideal, and regular local rings are domains. Divisors, the lemma that the blowup of an integral scheme is integral, and Algebra, the lemma that a regular local ring is a domain. Shows that the blowup is integral.
\item Characterization of regularity through closed points. Properties, the lemma characterizing regular schemes by regularity of the local rings at closed points. Reduces regularity to a check at the closed points of the blowup.
\item The exceptional divisor is an effective Cartier divisor isomorphic to the projective line. Divisors, the lemma that the blowup gives an effective Cartier divisor, and the identification of the exceptional divisor with the projective line over the residue field. Provides a local equation at each closed point whose quotient ring is the local ring of the projective line.
\item Regularity of the quotient by a regular element in a regular local ring. Algebra, the lemma deducing regularity of a quotient by a nonzerodivisor from regularity of the quotient, and the regularity of the local rings of affine space. Concludes regularity of the local ring of the blowup at a closed point from regularity of the local ring of the projective line.
\end{itemize}\medskip
\ifdefined\hypertarget\hypertarget{sourcepointer-1}{}\fi
\noindent\textit{Source reference.} Original lemma 00OP: sources/stacks-chow/tags.txt, source lines 903--903. Prerequisite pointer invoked at proof line(s) 750; it names a standard fact of a sibling chapter that the counted argument uses, and every step of the counted proof built on it is retained verbatim in the delivered body.\par\medskip
\ifdefined\hypertarget\hypertarget{sourcepointer-2}{}\fi
\noindent\textit{Source reference.} Original lemma 00NP: sources/stacks-chow/tags.txt, source lines 867--867. Prerequisite pointer invoked at proof line(s) 737; it names a standard fact of a sibling chapter that the counted argument uses, and every step of the counted proof built on it is retained verbatim in the delivered body.\par\medskip
\ifdefined\hypertarget\hypertarget{sourcepointer-3}{}\fi
\noindent\textit{Source reference.} Original lemma 00NU: sources/stacks-chow/tags.txt, source lines 872--872. Prerequisite pointer invoked at proof line(s) 752; it names a standard fact of a sibling chapter that the counted argument uses, and every step of the counted proof built on it is retained verbatim in the delivered body.\par\medskip
\ifdefined\hypertarget\hypertarget{sourcepointer-4}{}\fi
\noindent\textit{Source reference.} Original lemma 02ND: sources/stacks-chow/tags.txt, source lines 3476--3476. Prerequisite pointer invoked at proof line(s) 735; it names a standard fact of a sibling chapter that the counted argument uses, and every step of the counted proof built on it is retained verbatim in the delivered body.\par\medskip
\ifdefined\hypertarget\hypertarget{sourcepointer-5}{}\fi
\noindent\textit{Source reference.} Original lemma 02IT: sources/stacks-chow/tags.txt, source lines 3312--3312. Prerequisite pointer invoked at proof line(s) 740; it names a standard fact of a sibling chapter that the counted argument uses, and every step of the counted proof built on it is retained verbatim in the delivered body.\par\medskip
\begin{lemma}
\label{lemma-blowup-regular}
Let $(A, \mathfrak m, \kappa)$ be a regular local ring of dimension $2$.
Let $f : X \to S = \Spec(A)$ be the blowing up of $A$ in $\mathfrak m$.
Then $X$ is an irreducible regular scheme.
\end{lemma}

\begin{proof}
Observe that $X$ is integral by
Divisors, Lemma \ref{divisors-lemma-blow-up-integral-scheme}
and
Algebra, Lemma \ref{algebra-lemma-regular-domain}.
To see $X$ is regular it suffices to check that $\mathcal{O}_{X, x}$
is regular for closed points $x \in X$, see
Properties, Lemma \ref{properties-lemma-characterize-regular}.
Let $x \in X$ be a closed point. Since $f$ is proper $x$ maps to
$\mathfrak m$, i.e., $x$ is a point of the exceptional divisor $E$.
Then $E$ is an effective Cartier divisor and $E \cong \mathbf{P}^1_\kappa$.
Thus if $g \in \mathfrak m_x \subset \mathcal{O}_{X, x}$ is a local
equation for $E$, then
$\mathcal{O}_{X, x}/(g) \cong \mathcal{O}_{\mathbf{P}^1_\kappa, x}$.
Since $\mathbf{P}^1_\kappa$ is covered by two affine opens which are the
spectrum of a polynomial ring over $\kappa$, we see that
$\mathcal{O}_{\mathbf{P}^1_\kappa, x}$ is regular by
Algebra, Lemma \ref{algebra-lemma-dim-affine-space}.
We conclude by
Algebra, Lemma \ref{algebra-lemma-regular-mod-x}.
\end{proof}

\end{document}
